\bibitem{point-height-lan}
K.-W. Lan, \emph{Arithmetic compactifications of PEL-type Shimura
varieties}, revised thesis, 14 March 2021, Proposition 3.3.1.5
(extension of semiabelian homomorphisms) and Proposition 3.3.1.7
(closure of a generic subtorus), printed p.~192.
\url{https://www.kwlan.org/articles/cpt-PEL-type-thesis-revision.pdf}.
Both statements require the noetherian normal base retained here.

\bibitem{actual-point-height-source}
L. Alp\"oge, pinned source mirror
\path{jsmorph/rh}@\allowbreak
\texttt{ecd09ff9b4aaf6c08f2b5ef17dc49ca9648a77de}.
Under \path{Development/AutomorphicSpectrum/}:
\path{Definitions/MetricPeriods/IsAdapted.lean}, 11--65;
\path{Proofs/MetricProofs/LogIndexOfOrthogonalFamily.lean}, 563--636;
\path{Proofs/MetricProofs/ChartQContinuousAtAndLocalSolution.lean},
343--395 and 403--482; and
\path{Proofs/MetricProofs/PeriodDecompositionNearBoundary.lean},
62--170, 181--251, and 253--365.
The actual period construction is in
\path{Proofs/GeometryProofs/ComponentDataHodgeBound.lean}, 42--58,
\path{BoundaryChartLogPeriodFamily.lean}, 118--194, and
\path{SymplecticPolydiscPeriodBasis.lean}, 109--216 and 217--316.
The last two names are resolved to their complete paths in the
source certificate. The actual toric normalization is in
\path{Proofs/GeometryProofs/AbelianPartIntegralLattice.lean}, 37--86.
The determinant-degree tensor is
\path{Development/WeilBound/CoreStatements/PlurisubharmonicCones/}
\path{DetLatticeTensorIsDegOnMuTwDeg.lean}, 148--197.
The point identity and boundary comparison are in
\path{Proofs/GeometryProofs/CuspIntersectionHeight.lean},
232--315 and 350--426;
\path{CuspIntersectionComparison.lean}, 91--118 and 202--296;
and \path{CuspIntersectionFinite.lean}, 10--36.
The smoothing and the exact height consumer are
\path{Proofs/MetricProofs/RegularizedWeightDominatesEventually.lean},
78--107, and
\path{Proofs/GeometryProofs/CuspHeightUpper.lean}, 120--223.
The bounded source certificate and independent proof review
record hashes and distinguish these checks from a Lean rebuild.
